\chapter{Assignments for Independent Work}
\label{ch:tasks}

The assignments are arranged in order of increasing difficulty. The stars mark
the difficulty: \textcolor{FpgaOrange}{$\star$}---warm-up,
\textcolor{FpgaOrange}{$\star\star$}---you need to think,
\textcolor{FpgaOrange}{$\star\star\star$}---a mini-project. Each assignment
reinforces a particular chapter of Part~I (Table~\ref{tab:task-map}); the
final assignment 13 puts all the blocks together into the ``Security Panel''
and moves it onto the board.

\begin{table}[H]
\centering
\caption{Which chapter of Part~I each assignment reinforces}
\label{tab:task-map}
\small
\begin{tabularx}{\textwidth}{c X X}
\toprule
\textbf{No.} & \textbf{Assignment} & \textbf{Reinforces the chapter of Part~I} \\
\midrule
1  & Half adder on buttons & ``Arithmetic Blocks'' (Ch.~\ref{ch:arith}) \\
2  & Functions through NAND & ``Logic Gates and Boolean Algebra'' (Ch.~\ref{ch:gates}) \\
3  & Encoder on buttons & ``Decoders and Encoders'' (Ch.~\ref{ch:coders}) \\
4  & Binary counter on LEDs & ``Registers and Counters'' (Ch.~\ref{ch:regs}) \\
5  & Gray code counter & ``Number Systems and Codes'' (Ch.~\ref{ch:codes}) \\
6  & Running lights with buttons & ``Registers and Counters'' (Ch.~\ref{ch:regs}) \\
7  & Code lock & ``Synchronous Circuits, FSMs'' (Ch.~\ref{ch:sync}), ``FSM Design'' (Ch.~\ref{ch:fsmdesign}) \\
8  & Stopwatch & ``Registers and Counters'' (Ch.~\ref{ch:regs}), ``Synchronous Circuits'' (Ch.~\ref{ch:sync}) \\
9  & UART transmitter & ``Registers and Counters'' (Ch.~\ref{ch:regs}), ``FSM Design'' (Ch.~\ref{ch:fsmdesign}) \\
10 & UART receiver & ``Synchronous Circuits'' (Ch.~\ref{ch:sync}), the synchronizer \\
11 & The ``Reaction'' game & ``Registers and Counters'' (Ch.~\ref{ch:regs}): a shift register with feedback \\
12 & PWM for a servo & ``Digital and the Analog World'' (Ch.~\ref{ch:analog}) \\
13 & Security Panel on the board & all chapters of Part~I (the running project) \\
\bottomrule
\end{tabularx}
\end{table}

\begin{figure}[H]
\centering
\begin{tikzpicture}[node distance=5mm and 7mm,
  t/.style={block, minimum width=24mm, minimum height=11mm, font=\sffamily\scriptsize}]
  \node[t] (a) {1--3\\combinational\\logic};
  \node[t, right=of a] (b) {4--6\\counters\\and registers};
  \node[t, right=of b] (c) {7--8\\FSMs};
  \node[oblock, minimum width=24mm, minimum height=11mm, font=\sffamily\scriptsize, right=of c] (d) {9--10\\interfaces};
  \node[hblock, minimum width=24mm, minimum height=11mm, font=\sffamily\scriptsize, right=of d] (e) {11--13\\projects};
  \draw[arr] (a) -- (b); \draw[arr] (b) -- (c); \draw[arr] (c) -- (d); \draw[arr] (d) -- (e);
\end{tikzpicture}
\caption{Map of the assignments}
\end{figure}

\begin{taskbox}[Assignment 1 \quad \textcolor{FpgaOrange}{$\star$}]
Implement a \textbf{half adder} on two buttons: \code{led[0]} is the sum,
\code{led[1]} is the carry. Write the truth table and check it on the board.
\end{taskbox}

\begin{taskbox}[Assignment 2 \quad \textcolor{FpgaOrange}{$\star$}]
Express the AND, OR and NOT functions \textbf{using only NAND} and describe the
resulting circuit with the operators \code{\textasciitilde{}(... \& ...)}. How many
LUTs will the result take according to the synthesizer report? Explain why.
\end{taskbox}

\begin{taskbox}[Assignment 3 \quad \textcolor{FpgaOrange}{$\star\star$}]
Build an \textbf{encoder}: when one button is pressed the LEDs show the binary
code 1, for the other---2, for both---3. Use
\code{always @(*)} and \code{case}. Don't forget \code{default}!
\end{taskbox}

\begin{taskbox}[Assignment 4 \quad \textcolor{FpgaOrange}{$\star$}]
A binary counter on 6 LEDs that adds one 4 times per second.
After how many seconds will it overflow?
\end{taskbox}

\begin{taskbox}[Assignment 5 \quad \textcolor{FpgaOrange}{$\star\star$}]
A counter in \textbf{Gray code}: neighboring values differ in exactly one bit.
Hint: \code{gray = bin \^{} (bin >> 1)}. Compare the blinking with the
counter from assignment 4.
\end{taskbox}

\begin{taskbox}[Assignment 6 \quad \textcolor{FpgaOrange}{$\star\star$}]
Running lights in which button S1 changes the direction and S2 changes the
speed (3 speeds in a cycle).
\end{taskbox}

\begin{taskbox}[Assignment 7 \quad \textcolor{FpgaOrange}{$\star\star$}]
A \textbf{code lock}. The correct sequence of presses is:
S1, S1, S2, S1. After a correct entry all the LEDs light up; after
a mistake \code{led[5]} blinks and the entry starts over. Draw the
state transition graph of the machine.
\end{taskbox}

\begin{taskbox}[Assignment 8 \quad \textcolor{FpgaOrange}{$\star\star$}]
A \textbf{stopwatch}: S1 is start/stop, S2 is reset. The LEDs show the
number of seconds in binary (up to 63).
\end{taskbox}

\begin{taskbox}[Assignment 9 \quad \textcolor{FpgaOrange}{$\star\star\star$}]
A \textbf{UART transmitter}. Once per second send the character
\code{'A'} to the computer at 115\,200 baud (pin 69). The UART frame: a start bit (0),
8 data bits least significant bit first, a stop bit (1). The bit duration is
$27\,000\,000 / 115\,200 \approx 234$ clock cycles. Watch the result in any
terminal (PuTTY, \code{screen}, Arduino IDE).

\centering
\begin{tikztimingtable}[timing/slope=0.1, xscale=1.35, yscale=1.4, timing/font=\sffamily\small]
  TX & 2H 2L 2H 2L 2L 2L 2L 2L 2H 2L 2H 2H \\
\extracode
  \foreach \x/\t in {3/start, 5/D0, 7/D1, 9/D2, 11/D3, 13/D4, 15/D5, 17/D6, 19/D7, 21/stop}
    \node[font=\sffamily\tiny] at (\x,1.75) {\t};
\end{tikztimingtable}
\par\small Transmission of the character 'A' = \code{0x41} = \code{0100\,0001}$_2$
\end{taskbox}

\begin{taskbox}[Assignment 10 \quad \textcolor{FpgaOrange}{$\star\star\star$}]
A \textbf{UART receiver}: show the byte received from the computer on the
LEDs (the low 6 bits). The input is pin 70. Don't forget the synchronizer!
\end{taskbox}

\begin{taskbox}[Assignment 11 \quad \textcolor{FpgaOrange}{$\star\star\star$}]
The \textbf{``Reaction'' game}. After a random time (use a
linear-feedback shift register, LFSR) an LED lights up.
The player must press S1 as fast as possible. The reaction time in tens of
milliseconds is shown on the LEDs or sent over UART.
\end{taskbox}

\begin{taskbox}[Assignment 12 \quad \textcolor{FpgaOrange}{$\star\star\star$}]
A \textbf{PWM generator for a servo}. Servos and drone motor controllers
are driven by pulses 1--2~ms long with a period of 20~ms.
Build a generator in which buttons S1 and S2 change the pulse width in
steps of 0.1~ms, and output the signal to any free GPIO pin. Check it
with an oscilloscope or with a real servo.
\end{taskbox}

\begin{taskbox}[Assignment 13 \quad \textcolor{FpgaOrange}{$\star\star\star$}\quad ``Security Panel'' on the board]
\textbf{Final assignment.} Build on the Tang Nano 20K the security panel that you developed in Part~I and described in Verilog in Chapter~\ref{ch:bridge}: three sensors (door, window, motion), night mode, a noise threshold, arming, alarm memory, entering the code 1-0-1, a siren (an LED or a buzzer). The board has only two buttons, S1 and S2, so move some of the inputs to free GPIO pins with toggle switches (or build the circuit on the ``Lab'' tab): S1 is the door, S2 is arming and entering a code bit, the rest go on toggle switches.

\textbf{How to proceed.}
\begin{enumerate}
  \item In Logra create a ``From scratch'' project, enter the \code{guard\_panel} module from Chapter~\ref{ch:bridge}, press ``Synthesize (Ctrl+Enter)'' and check the behavior on the ``Waveforms'' tab.
  \item Add a synchronizer and debouncing to each button, and a synchronous reset (Chapter~\ref{ch:clocks}).
  \item ``Pins'' tab: assign the buttons, LEDs and the buzzer; save the project. ``Chip'' tab: ``Place and build,'' $f_{max}$ not lower than 54~MHz.
  \item ``Connect'' tab: load the design into SRAM and check it on the board.
\end{enumerate}

\textbf{Acceptance criterion.}
\begin{enumerate}
  \item The panel is armed by the button and does not react to the sensors in the OFF state.
  \item An alarm is triggered by any of three causes: two sensors out of three, noise above the threshold, or the door at night.
  \item The alarm is remembered and does not go off when the door is closed.
  \item The alarm can be cleared only by entering the code 1-0-1; any other sequence does not work.
  \item When a button bounces, the state changes exactly once; $f_{max}$ in the nextpnr report is not lower than 54~MHz.
\end{enumerate}
\end{taskbox}
